% ------------------------------------------------------------------
%  Capitolo 21 — Ogni disciplina ha il suo studio
%  Parte V
%  Ricerca di riferimento: ricerca DA RICERCARE (ricerca 07, discipline universitarie); 01 §5.3, §16
%  Voci bibliografiche principali: murray2025, rohrer2007, bahrick1993
%  Epigrafe: ◐ verificata indirettamente (vedi ricerca 03 e registro, ricerca 00)
%  Scheda del capitolo: ricerca/06_indice-ragionato.md, capitolo 21
% ------------------------------------------------------------------
\chapter{Ogni disciplina ha il suo studio}
\label{cap:discipline}

\epigrafe[le cose che bisogna aver imparato per farle, le impariamo facendole]{ἃ γὰρ δεῖ μαθόντας ποιεῖν, ταῦτα ποιοῦντες μανθάνομεν}{Aristotele, Etica Nicomachea II 1, 1103a32}

\iniziale{Q}{uesto capitolo} \segnaposto{apertura: da scrivere — vedi la scheda nell'indice ragionato}.

\section{Matematica, fisica e ingegneria}

\segnaposto{da scrivere}

\section{Scienze della vita, medicina e professioni sanitarie}

\segnaposto{da scrivere}

\section{Giurisprudenza ed economia}

\segnaposto{da scrivere}

\section{Discipline umanistiche}

\segnaposto{da scrivere}

\section{Lingue straniere e classiche}

\segnaposto{da scrivere}

\begin{daricordare}
\segnaposto{il principio del capitolo in due righe}
\end{daricordare}

\begin{domande}
  \item \segnaposto{domanda di recupero su questo capitolo}
  \item \segnaposto{domanda di applicazione a una materia che studi}
  \item \segnaposto{domanda di ripresa su un capitolo precedente}
\end{domande}

\begin{sintesi}
  \item \segnaposto{punto di sintesi}
\end{sintesi}

\finecapitolo
